% ============================================================================
% LANGENTWURF LE-07a: Comida (Lebensmittel)
% Spanisch Klasse 7 - Sequenz 7: Compras (Einkaufen)
% ============================================================================
% Tier: 1 (vollständig)
% Doppelstunde: 1-2 von 8
% Fachfarbe: #c41e3a (Karminrot)
% ============================================================================

\documentclass[11pt,a4paper]{article}

\usepackage[utf8]{inputenc}
\usepackage[T1]{fontenc}
\usepackage[ngerman]{babel}
\usepackage{geometry}
\usepackage{fancyhdr}
\usepackage{lastpage}
\usepackage{xcolor}
\usepackage{tcolorbox}
\usepackage{enumitem}
\usepackage{tabularx}
\usepackage{longtable}
\usepackage{array}
\usepackage{booktabs}
\usepackage{amssymb}

% ============================================================================
% KONFIGURATION
% ============================================================================

\newcommand{\fach}{Spanisch}
\newcommand{\fachfarbe}{c41e3a}
\newcommand{\jahrgangsstufe}{7}
\newcommand{\schulart}{Gesamtschule}
\newcommand{\stundenthema}{Comida -- Lebensmittel}
\newcommand{\reihenthema}{Compras (Einkaufen)}
\newcommand{\stundennummer}{7a}
\newcommand{\reihenstunden}{8}
\newcommand{\gerniveau}{A1}
\newcommand{\lehrwerk}{Qué pasa 1}
\newcommand{\lehrwerkseite}{S. 124-128}
\newcommand{\lerngruppengroesse}{24}

% ============================================================================
% LAYOUT
% ============================================================================
\geometry{left=2.5cm, right=2.5cm, top=2.5cm, bottom=2.5cm}
\definecolor{fach}{HTML}{\fachfarbe}
\definecolor{gray}{HTML}{666666}
\definecolor{lightgray}{HTML}{f5f5f5}
\definecolor{successgreen}{HTML}{2a7a4b}
\definecolor{warningorange}{HTML}{b35c00}
\definecolor{basis}{HTML}{2a7a4b}
\definecolor{standard}{HTML}{2c5aa0}
\definecolor{erweiterung}{HTML}{7b1fa2}

\tcbuselibrary{skins,breakable}

\newtcolorbox{infobox}[1][]{
    colback=lightgray,
    colframe=fach,
    fonttitle=\bfseries,
    title=#1,
    breakable
}

\newtcolorbox{scorebox}{
    colback=white,
    colframe=fach,
    fonttitle=\bfseries,
    title=Didaktik-Score,
    breakable
}

\pagestyle{fancy}
\fancyhf{}
\fancyhead[L]{\small\textcolor{gray}{\fach{} -- Jg. \jahrgangsstufe{} -- \stundenthema}}
\fancyhead[R]{\small\textcolor{gray}{LE-\stundennummer{} | Tier 1}}
\fancyfoot[C]{\small\thepage{} / \pageref{LastPage}}
\renewcommand{\headrulewidth}{0.4pt}

% Differenzierungsmarker
\newcommand{\B}{\textcolor{basis}{\textbf{[B]}}}
\newcommand{\Std}{\textcolor{standard}{\textbf{[S]}}}
\newcommand{\E}{\textcolor{erweiterung}{\textbf{[E]}}}

% ============================================================================
% DOKUMENT
% ============================================================================
\begin{document}

% ============================================================================
% DECKBLATT
% ============================================================================
\begin{titlepage}
    \centering
    \vspace*{2cm}
    {\large\textcolor{gray}{\schulart}}\\[2cm]
    {\Huge\textcolor{fach}{\textbf{Langentwurf}}}\\[0.5cm]
    {\large LE-\stundennummer{} | Tier 1 (vollständig)}\\[2cm]
    {\LARGE\textbf{\stundenthema}}\\[0.5cm]
    {\large Sequenz: \reihenthema}\\[2cm]
    \begin{tabular}{rl}
        \textbf{Fach:} & \fach{} (\gerniveau) \\[0.3cm]
        \textbf{Klasse:} & \jahrgangsstufe \\[0.3cm]
        \textbf{Lehrwerk:} & \lehrwerk, \lehrwerkseite \\[0.3cm]
        \textbf{Doppelstunde:} & \stundennummer{} (Std. 1-2 von 8) \\[0.3cm]
    \end{tabular}
    \vfill
\end{titlepage}

\tableofcontents
\newpage

% ============================================================================
% 1. BEDINGUNGSANALYSE
% ============================================================================
\section{Bedingungsanalyse}

\subsection{Lerngruppe}
\begin{tabular}{ll}
    Kursgröße: & \lerngruppengroesse{} SuS \\
    Jahrgangsstufe: & \jahrgangsstufe \\
    Sprachniveau: & \gerniveau{} (GER) -- Anfänger \\
    Fremdsprache: & 2. Fremdsprache (nach Englisch) \\
\end{tabular}

\subsection{Vorwissen aus Sequenz 1-6}
\begin{itemize}
    \item Begrüßung, Vorstellung, Verabschiedung
    \item Verb \textit{ser}, \textit{tener}, \textit{hay}
    \item Regelmäßige Verben auf -ar, -er, -ir
    \item Artikel (el, la, los, las, un, una)
    \item Grundzahlen 1-100
    \item Farben, Familie, Schule, Stadt
\end{itemize}

\subsection{Differenzierung (intern)}
\begin{tabular}{|l|p{3.5cm}|p{3.5cm}|p{3.5cm}|}
\hline
& \B{} \textbf{Basis} & \Std{} \textbf{Standard} & \E{} \textbf{Erweiterung} \\
\hline
\textbf{Ziel} & BR: Rezeption, Benennung & MR: Kategorisierung, Anwendung & GY: Produktion, Transfer \\
\hline
\textbf{Vokabeln} & 10 Kernwörter mit Bildunterstützung & 15 Wörter + Kategorien & 20+ Wörter + Mengenangaben \\
\hline
\textbf{Aufgaben} & Zuordnung Bild-Wort & Kategorisieren, Lückentext & Einkaufsliste schreiben \\
\hline
\end{tabular}

\vspace{0.5em}
\textbf{WICHTIG:} Diese Marker sind NUR für die Lehrkraft!

% ============================================================================
% 2. SACHANALYSE
% ============================================================================
\section{Sachanalyse}

\subsection{Wortfeld Lebensmittel (Comida)}

\textbf{Obst (La fruta):}
\begin{itemize}
    \item la manzana (der Apfel)
    \item el plátano (die Banane)
    \item la naranja (die Orange)
    \item la fresa (die Erdbeere)
    \item la uva (die Traube)
\end{itemize}

\textbf{Gemüse (La verdura):}
\begin{itemize}
    \item el tomate (die Tomate)
    \item la lechuga (der Salat)
    \item la zanahoria (die Karotte)
    \item la cebolla (die Zwiebel)
    \item el pimiento (die Paprika)
\end{itemize}

\textbf{Fleisch/Fisch (La carne / El pescado):}
\begin{itemize}
    \item la carne (das Fleisch)
    \item el pollo (das Hähnchen)
    \item el pescado (der Fisch)
\end{itemize}

\textbf{Grundnahrungsmittel:}
\begin{itemize}
    \item el pan (das Brot)
    \item la leche (die Milch)
    \item el queso (der Käse)
    \item el huevo (das Ei)
    \item el arroz (der Reis)
\end{itemize}

\subsection{Mengenangaben}

\begin{center}
\begin{tabular}{|l|l|}
\hline
\textbf{Spanisch} & \textbf{Deutsch} \\
\hline
un kilo de... & ein Kilo... \\
medio kilo de... & ein halbes Kilo... \\
un litro de... & ein Liter... \\
una botella de... & eine Flasche... \\
un paquete de... & ein Paket/eine Packung... \\
una docena de... & ein Dutzend... \\
\hline
\end{tabular}
\end{center}

\subsection{Kontrastive Betrachtung}

\begin{tabular}{|l|p{5cm}|p{5cm}|}
\hline
& \textbf{Deutsch} & \textbf{Spanisch} \\
\hline
Genus & Tomate = feminin & tomate = maskulin (el tomate) \\
\hline
Artikel & Partitiv mit ``ein'' & Kein Partitiv: Quiero pan (nicht: *un pan) \\
\hline
Pluralbildung & Äpfel & manzanas (regelmäßig) \\
\hline
\end{tabular}

\textbf{Typische Interferenz:} *``la tomate'' (falscher Artikel durch dt. Genus)

% ============================================================================
% 3. DIDAKTISCHE ANALYSE
% ============================================================================
\section{Didaktische Analyse}

\subsection{Stellung in der Sequenz}
\begin{tabular}{|c|l|l|}
\hline
\textbf{Block} & \textbf{Thema} & \textbf{Status} \\
\hline
\textbf{7a} & \textbf{Comida (Lebensmittel)} & \textbf{diese Stunde} \\
7b & En el supermercado (Dialog) & folgt \\
7c & Preise und Zahlen & folgt \\
7d & Einkaufen -- Sequenztest & folgt \\
\hline
\end{tabular}

\subsection{Kompetenzziele (Rahmenplan MV)}
\begin{itemize}
    \item \textbf{Kommunikativ:} Sprechen (Benennung), Hörverstehen
    \item \textbf{Sprachlich:} Wortschatz (Lebensmittel), Grammatik (Mengenangaben mit \textit{de})
    \item \textbf{Interkulturell:} Spanische Esskultur, Märkte
    \item \textbf{Methodisch:} Wortschatzarbeit, Kategorisierung, Memory
\end{itemize}

\subsection{Legitimation}
\textbf{Rahmenplan Spanisch MV (Kl. 7-10):}
\begin{itemize}
    \item Kompetenzbereich: Verfügung über sprachliche Mittel -- Wortschatz
    \item Standard: ``Die SuS können Grundbegriffe aus dem Bereich Lebensmittel und Einkaufen benennen und anwenden.''
    \item Themenfeld: Alltag -- Essen und Trinken
\end{itemize}

% ============================================================================
% 4. STUNDENZIELE
% ============================================================================
\section{Stundenziele}

\subsection{Grobziel}
Die SuS erweitern ihre \textbf{Wortschatzkompetenz im Bereich Lebensmittel}, indem sie Vokabeln aus verschiedenen Kategorien (Obst, Gemüse, Fleisch, Grundnahrungsmittel) \textbf{benennen}, \textbf{kategorisieren} und mit Mengenangaben \textbf{verwenden}.

\subsection{Feinziele (operationalisiert)}
\begin{tabular}{|c|p{7.5cm}|c|c|}
\hline
\textbf{Nr.} & \textbf{Die SuS können...} & \textbf{AFB} & \textbf{Diff.} \\
\hline
FZ1 & 15 Lebensmittelbegriffe aus den Kategorien Obst, Gemüse, Fleisch/Fisch und Grundnahrungsmittel \textbf{benennen} und den Bildern \textbf{zuordnen} & I & alle \\
\hline
FZ2 & Lebensmittel nach Kategorien (fruta, verdura, carne) \textbf{ordnen} und die Oberbegriffe \textbf{anwenden} & I/II & \Std{}/\E{} \\
\hline
FZ3 & Mengenangaben (un kilo de, medio kilo de, un litro de) korrekt mit Lebensmitteln \textbf{verbinden} und in Sätzen \textbf{verwenden} & II & \Std{}/\E{} \\
\hline
FZ4 & eine einfache Einkaufsliste auf Spanisch \textbf{erstellen} und \textbf{präsentieren} & III & \E{} \\
\hline
\end{tabular}

\vspace{1em}
\textbf{AFB-Verteilung:} AFB I: 30\% | AFB II: 45\% | AFB III: 25\%

% ============================================================================
% 5. METHODISCHE ANALYSE
% ============================================================================
\section{Methodische Analyse}

\subsection{Einstieg}
\begin{itemize}
    \item \textbf{Methode:} Bildimpuls -- Foto eines spanischen Marktes (Mercado de la Boquería)
    \item \textbf{Begründung:} Visuelle Aktivierung, kultureller Bezug, Neugier wecken
    \item \textbf{Alternative:} Realia (echte Lebensmittel / Plastikobst)
\end{itemize}

\subsection{Erarbeitung}
\begin{itemize}
    \item \textbf{Methode:} Wortschatzeinführung mit Bild-Wort-Karten
    \item \textbf{Begründung:} Visuelle Unterstützung für Vokabellernen
    \item \textbf{Scaffolding:} Kategorisierung nach Obst/Gemüse/etc.
\end{itemize}

\subsection{Übungsphasen}
\begin{itemize}
    \item \textbf{Methode:} Memory-Spiel $\rightarrow$ Kategorisierung $\rightarrow$ Lückentext $\rightarrow$ Einkaufsliste
    \item \textbf{Progression:} Rezeptiv $\rightarrow$ Reproduktiv $\rightarrow$ Produktiv
    \item \textbf{Differenzierung:} Anzahl der Vokabeln, Hilfestellung
\end{itemize}

\subsection{Medien}
\begin{itemize}
    \item Tafel/Whiteboard: Kategorien-Tabelle
    \item Lehrwerk: \lehrwerk, \lehrwerkseite
    \item Arbeitsblatt: AB-07a.pdf
    \item Lernpfad: LP-07a.html (Tablets, optional)
    \item Bild-Wort-Karten (Memory-Set)
    \item Bildmaterial: Marktfotos
\end{itemize}

% ============================================================================
% 6. VERLAUFSPLANUNG
% ============================================================================
\section{Verlaufsplanung}

\textbf{1. Stunde (45 Min.)}

\begin{longtable}{|p{1cm}|p{1.5cm}|p{4cm}|p{3cm}|p{1.5cm}|p{2cm}|}
\hline
\textbf{Zeit} & \textbf{Phase} & \textbf{Lehrerhandeln} & \textbf{Schülerhandeln} & \textbf{SF} & \textbf{Material} \\
\hline
\endhead

0--5 & Einstieg & Bild Mercado zeigen: ``¿Qué veis?'' & Vermutungen, bekannte Wörter nennen & Plenum & Bild \\
\hline

5--15 & Input 1 & Vokabeln Obst (fruta) einführen mit Bildern: manzana, plátano, naranja & Nachsprechen, Artikelzuordnung & Plenum & Karten \\
\hline

15--25 & Input 2 & Vokabeln Gemüse (verdura): tomate, lechuga, zanahoria & Nachsprechen, Eintrag AB & Plenum & Karten, AB \\
\hline

25--35 & Übung 1 & Memory-Spiel in Gruppen (Bild $\leftrightarrow$ Wort) & Spielen, Vokabeln wiederholen & GA & Memory \\
\hline

35--40 & Sicherung & Schnelles Quiz: ``¿Cómo se dice Apfel?'' & Antworten & Plenum & -- \\
\hline

40--45 & Ausblick & Vorschau Grundnahrungsmittel, HA & Notieren & Plenum & -- \\
\hline
\end{longtable}

\vspace{1em}

\textbf{2. Stunde (45 Min.)}

\begin{longtable}{|p{1cm}|p{1.5cm}|p{4cm}|p{3cm}|p{1.5cm}|p{2cm}|}
\hline
\textbf{Zeit} & \textbf{Phase} & \textbf{Lehrerhandeln} & \textbf{Schülerhandeln} & \textbf{SF} & \textbf{Material} \\
\hline
\endhead

0--5 & Aktivier. & Wiederholung: Obst/Gemüse benennen & Antworten & Plenum & Karten \\
\hline

5--15 & Input 3 & Fleisch/Fisch + Grundnahrungsmittel: carne, pollo, pescado, pan, leche, queso & Nachsprechen, Eintrag AB & Plenum & Karten, AB \\
\hline

15--25 & Input 4 & Mengenangaben einführen: un kilo de..., medio kilo de..., un litro de... & Beispiele nachsprechen & Plenum & Tafel \\
\hline

25--35 & Übung 2 & Lückentext AB: Mengenangaben einsetzen & EA: AB bearbeiten & EA & AB-07a \\
\hline

35--40 & Anwendung & \E{}: Einkaufsliste schreiben (5 Produkte mit Mengen) & Schreiben, präsentieren & EA/PA & AB-07a \\
\hline

40--45 & Abschluss & Zusammenfassung, HA: Vokabeln lernen & Notieren & Plenum & -- \\
\hline
\end{longtable}

\textbf{Legende:} EA = Einzelarbeit, PA = Partnerarbeit, GA = Gruppenarbeit

% ============================================================================
% 7. ERWARTETE SCHWIERIGKEITEN
% ============================================================================
\section{Erwartete Schwierigkeiten}

\begin{tabular}{|p{4.5cm}|p{8cm}|}
\hline
\textbf{Schwierigkeit} & \textbf{Reaktion} \\
\hline
Genus \textit{el tomate} (nicht la) & Expliziter Hinweis, kontrastiv zum Deutschen \\
\hline
Aussprache \textit{zanahoria} & Langsam vorsprechen, Silbenweise: za-na-ho-ria \\
\hline
Verwechslung \textit{manzana/naranja} & Visuelle Hilfe (Farbe), Eselsbrücken \\
\hline
\textit{de} bei Mengenangaben vergessen & Struktur explizit: \textit{un kilo DE manzanas} \\
\hline
Artikelverwendung bei Mengen & Im Spanischen: medio kilo de pan (ohne Artikel) \\
\hline
Zeitproblem & Memory kürzen, Einkaufsliste als HA \\
\hline
\end{tabular}

% ============================================================================
% 8. HAUSAUFGABE
% ============================================================================
\section{Hausaufgabe}

\begin{itemize}
    \item Lehrwerk S. 126, Übung 2 (Vokabeln zuordnen)
    \item Vokabeln lernen: 15 Lebensmittel + 5 Mengenangaben
    \item Optional: Lernpfad LP-07a.html bearbeiten
    \item \E{}: Eigene Einkaufsliste für ein Wochenende schreiben (10 Produkte)
\end{itemize}

% ============================================================================
% 9. LÖSUNGEN
% ============================================================================
\section{Lösungen}

\subsection{AB-07a Lösungen}

\textbf{Aufgabe 1: Zuordnung Bild-Wort}
\begin{itemize}
    \item Apfel $\rightarrow$ la manzana
    \item Banane $\rightarrow$ el plátano
    \item Orange $\rightarrow$ la naranja
    \item Tomate $\rightarrow$ el tomate
    \item Salat $\rightarrow$ la lechuga
    \item Hähnchen $\rightarrow$ el pollo
    \item Fisch $\rightarrow$ el pescado
    \item Brot $\rightarrow$ el pan
    \item Milch $\rightarrow$ la leche
    \item Käse $\rightarrow$ el queso
\end{itemize}

\textbf{Aufgabe 2: Kategorisierung}
\begin{itemize}
    \item \textbf{La fruta:} manzana, plátano, naranja
    \item \textbf{La verdura:} tomate, lechuga, zanahoria
    \item \textbf{La carne/El pescado:} pollo, pescado, carne
    \item \textbf{Otros:} pan, leche, queso
\end{itemize}

\textbf{Aufgabe 3: Lückentext Mengenangaben}
\begin{enumerate}[label=\alph*)]
    \item Quiero \textbf{un kilo de} manzanas.
    \item Necesito \textbf{medio kilo de} tomates.
    \item Dame \textbf{un litro de} leche.
    \item Quiero \textbf{una botella de} agua.
\end{enumerate}

\textbf{Aufgabe 4: Einkaufsliste (Beispiel)}
\begin{itemize}
    \item un kilo de manzanas
    \item medio kilo de tomates
    \item un litro de leche
    \item un paquete de arroz
    \item 200 gramos de queso
\end{itemize}

% ============================================================================
% 10. ANLAGEN
% ============================================================================
\section{Anlagen}

\begin{enumerate}
    \item Arbeitsblatt (AB-07a.pdf)
    \item Musterlösung (ML-07a.pdf)
    \item Lehrerhinweise (LH-07a.pdf)
    \item Lernpfad (LP-07a.html)
    \item Bild-Wort-Karten (Memory-Set, Kopiervorlage)
    \item Marktfoto (Mercado de la Boquería)
\end{enumerate}

% ============================================================================
% 11. DIDAKTIK-SCORE
% ============================================================================
\newpage
\section{Didaktik-Score}

\begin{scorebox}
\textbf{Bewertung der didaktischen Qualität nach 10 Dimensionen}\\
\textbf{Ziel-Score: $\geq$ 9.0 (Premium-Qualität)}

\vspace{1em}
\begin{tabular}{|c|l|c|c|p{4cm}|}
\hline
\textbf{\#} & \textbf{Dimension} & \textbf{Gew.} & \textbf{Score} & \textbf{Notiz} \\
\hline
1 & Differenzierung & 15\% & 9/10 & [B]/[S]/[E] qualitativ verschieden \\
\hline
2 & Taxonomie (AFB) & 10\% & 9/10 & 30/45/25 gut verteilt \\
\hline
3 & Lernziel-Alignment & 10\% & 10/10 & RP MV Wortschatz explizit \\
\hline
4 & Aktivierung & 10\% & 9/10 & Memory, Kategorisierung aktiv \\
\hline
5 & Scaffolding & 10\% & 9/10 & Bild-Wort, Kategorien \\
\hline
6 & Feedback-Qualität & 10\% & 9/10 & LP: elaboriert \\
\hline
7 & Sprachliche Klarheit & 10\% & 10/10 & Kl. 7 angemessen \\
\hline
8 & Motivation & 10\% & 10/10 & Alltagsbezug Essen, Markt \\
\hline
9 & Inklusion (UDL) & 10\% & 9/10 & Bilder, Karten, Kategorien \\
\hline
10 & Praktikabilität & 5\% & 9/10 & 90 Min. realistisch \\
\hline
\multicolumn{3}{|r|}{\textbf{GESAMT:}} & \textbf{9.3/10} & \\
\hline
\end{tabular}

\vspace{1em}
$\checkmark$ \textcolor{successgreen}{\textbf{FREIGABE} (Score $\geq$ 9.0)}
\end{scorebox}

\end{document}
